% !TEX program = pdflatex
\documentclass[11pt]{article}
\usepackage{support-notes}
\title{A lower bound for derived tensor products from a double derivation}
\author{\workingauthor}\date{\coursedate}
\begin{document}\maketitle
\begin{abstract}
A cycle admitting a noncommutative partial derivative recovers three
independent homology classes from their product with two inserted copies
of that cycle. For perfect modules this gives a lower bound in each
degree, including after passage to a dg algebra quotient. In particular
over a positive free product every tensor of perfect modules with two
infinite-cohomology factors contains, degree by degree, a shifted copy of
the algebra's cohomology. All theorem statements below are proved.
\end{abstract}

\section{The estimate}
Fix a field $k$, cochain conventions and a morphism $\pi:A\to G$ of unital
dg algebras concentrated in nonnegative degrees. Assume their graded pieces
are finite dimensional. A finite cell module has underlying graded module
a finite sum of shifted copies of the algebra, with a semi-free
filtration. Perfect means a retract in the derived category of such a
module. Let $E$ be perfect right and $Y$ perfect left over $A$.
Put $\bar E=E\ten_A G$ and $\bar Y=G\ten_A Y$. For a complex $M$ let
$H^iM$ be its $i$th cohomology. Write
\[
 e_p=\rank(H^pE\to H^p\bar E),\quad
 g_m=\rank(H^m A\to H^m G),\quad
 y_q=\rank(H^qY\to H^q\bar Y).
\]
The outer maps are induced by $\pi$ on underlying complexes.
Take a one-dimensional complex $L=k\ell$ with $|\ell|=j>0$ and zero
differential. Suppose there is a degree-zero chain map
$\delta:G\to G\otimes L\otimes G$ with
\[
 \delta(bc)=\delta(b)c+b\delta(c),\qquad
 \delta(\pi v)=1\otimes\ell\otimes1
\]
for a closed $v\in A^j$. The target has the ordinary outer dg bimodule
structure. Placing the degree in the middle slot avoids sign ambiguities
in the Leibniz rule. A basic instance is a free product $G=C*T(kv)$,
with $dv=0$, using the derivation that kills $C$.

\begin{thm}\label{cut}
There are bounds $p_0(E),q_0(Y)$ such that, for $p\ge p_0$,
$q\ge q_0$ and every $m\ge0$,
\[
          \dim_k H^{p+j+m+j+q}(E\ten_A Y)\ \ge\ e_p g_m y_q .
\]
More precisely the kernel of insertion $[e]\otimes[a]\otimes[y]
\mapsto[e v a v\otimes y]$ in these degrees is contained in the kernel
of the tensor product of the three reduction maps. The one-insertion
map $[e]\otimes[y]\mapsto[e v\otimes y]$ gives the analogous bound
$e_p y_q$ in degree $p+j+q$.
\end{thm}

\begin{proof}
First take finite cell representatives with graded bases $e_i,f_b$.
Their ordinary tensor computes the derived one; its vectors have
coordinates $e_i a f_b$. Define a $k$-linear two-cut map on this tensor
by applying $\delta^{(2)}=(1\otimes1\otimes\delta)\delta$ to $\pi(a)$
and placing the outer factors in $\bar E$ and $\bar Y$ using their
indicated basis vectors. Project the three non-marker factors to the
cochain quotients
\[
 \tau^{\ge p}\bar E\ \otimes L\otimes\
 \tau^{\ge m}G\ \otimes L\otimes\ \tau^{\ge q}\bar Y.           \tag{1}
\]
These are truncations of $k$-complexes: degree $p$ in the first is
$\bar E^p/d\bar E^{p-1}$ and higher degrees are untouched. No module
action is needed on a truncation.

This projected map commutes with differentials. For clarity, before
projection the two-cut product rule is the sum of cuts both in the
first factor, a cut in each factor, and both cuts in the second factor;
in the middle term one multiplies the two adjacent unmarked slots.
It follows by applying the degree-zero derivation twice. The internal
differential on $a$ commutes with both cuts. In a term
$d e_i=\sum_h e_h c_{hi}$ the no-cut contribution in $c_{hi}$ is precisely
the differential of the outer basis on the target. Any contribution
cutting $c_{hi}$ has left degree at most
$|e_h|+|c_{hi}|-j=|e_i|+1-j$ because all the remaining algebra degrees
are nonnegative. Choose $p$ strictly greater than this for every $i$.
Similarly a cut in a coefficient of $df_b$ has right degree at most
$|f_b|+1-j<q$. This accounts for every extra contribution. The signs of
the no-cut terms agree because the displayed maps on coefficients have
total degree zero and retain the degrees of both markers; for example the
sign of the right-basis differential is $(-1)^{|e_i|+|a|}$ on either
side. The derivations are chain maps, so these observations establish the
claim even with nonzero internal differential on $G$.

Apply this chain map to a product of cycles $e v a v y$ with outer
degrees $p,q$ and middle degree $m$. A cut inside $e$ leaves degree
at most $p-j$ on the left, and a cut inside $y$ leaves degree at most
$q-j$ on the right. Of the remaining terms only the cuts at the two
specified $v$'s can survive (1). A first cut later than the first marker
or a last cut earlier than the last marker consumes strictly positive
degree from the middle slot; the term is killed by one of the three
truncations. More explicitly, the sum of the three output degrees is
$p+m+q$, hence their three lower bounds must all be equalities.
Cutting a marker twice gives zero because its derivative has scalar
outer factors. The surviving tensor is precisely
$\bar e\otimes\ell\otimes\pi(a)\otimes\ell\otimes\bar y$ after projection.
The Künneth isomorphism over $k$ and preservation of cohomology at each
cutoff identify its rank as $e_p g_m y_q$. Insertion itself is a chain
map from the tensor of complexes with the two degree-$j$ markers,
because $dv=0$. The same argument with a single derivative and the two
outer truncations proves the one-cut assertion.

For retracts choose cofibrant representatives and split inclusions in
the derived category into finite cell modules. Derived reduction and
tensor preserve those retracts. Insertion is natural in the outer
module variables: on cofibrant models this is ordinary right and left
linearity, and passing to homology respects homotopies. Therefore a
relation between insertions for retracts maps to one between insertions
in the finite cell representatives. Their reductions on homology are
split injections on the retracts. The finite cell bounds suffice.
\end{proof}

\begin{cor}\label{rigid}
If $\pi$ is the identity and $H^*G$ is infinite, then
$H^*(E\ten_G Y)$ is finite dimensional if and only if $H^*E$ or $H^*Y$
is finite dimensional. If both are infinite there is an integer $b$ such
that $\dim H^{m+b}(E\ten_G Y)\ge\dim H^mG$ for every $m\ge0$.
\end{cor}
\begin{proof}
Choose nonzero $H^p E,H^qY$ above the bounds. Local finiteness and a
lower degree bound ensure such degrees exist. Apply Theorem~\ref{cut}.
For the converse, tensor the finite-cohomology factor with successive
cells, shifts and summands of the perfect factor on the other side.
\end{proof}

In particular the exponential growth rate $\limsup_d(\dim H^d)^{1/d}$
of such an unbounded tensor equals that of $G$ when the latter exceeds
one: the upper bound follows by finite triangles from finitely many
shifts of $G$. Substituting $Y=P^\vee=\RH_G(P,G)$ gives the Ext version.
It makes every profile either the full category or $\UU_G$, and both
occur. Noncommutativity, nonsplit residue extensions and characteristic
two require no alteration.

\paragraph{Proved failure without a derivative.}
In $G=k[x,z]$ with $|x|=|z|=2$ and zero differential, $G/(x)$ and
$G/(z)$ are perfect and both have infinite cohomology, while their
derived tensor has only $k$ in degree zero. The candidate derivative
$\delta(z)=1\otimes\ell\otimes1,\delta(x)=0$ violates $xz=zx$.
A derivation to ordinary commutative differentials cannot be substituted
in the proof; the two independent outer slots are used.

\paragraph{Proved relative version for a normal element.}
Suppose $A$ has zero differential, $h$ is a positive homogeneous regular
normal element, $G=A/(h)$, and the above derivative and lift exist.
If right multiplication by $h$ on $H^*E$ and left multiplication by $h$
on $H^*Y$ both fail to be eventually bijective, the lower bound of
Corollary~\ref{rigid} holds with middle $H^mG$. Indeed the respective
images in reduced cohomology are the cokernels of multiplication by
$h$ in each degree. For the right factor use the short exact sequence
of left modules $0\to Ah\to A\to G\to0$ and the free left generator
$h$ of the ideal. For the left factor use the right generator of $hA$.
These arguments concern underlying complexes and include the usual shift
sign in the left-handed case. In particular they do not require $h$
to be in the graded centre. Eventual surjectivity implies eventual
bijectivity on locally finite cohomology: along each residue class modulo
$|h|$ the dimensions eventually decrease weakly, hence stabilize.
Thus arbitrarily high $e_p,y_q$ are positive, and $H(A)\to H(G)$ is onto.
This argument avoids replacing an odd central element by an untwisted
dg transformation.

\paragraph{Literature and status.}
Outer double derivations are standard in
\href{https://arxiv.org/abs/math/0410528}{Van den Bergh}.
Word resolutions over fibre products occur in
\href{https://arxiv.org/abs/0704.3631}{Moore}. Rigidity of eventual
Tor over commutative dg trivial extensions and over ordinary fibre
products is proved by
\href{https://arxiv.org/abs/1508.00748}{Avramov--Iyengar--Nasseh--Sather-Wagstaff}
and \href{https://arxiv.org/abs/1603.05711}{Nasseh--Sather-Wagstaff}.
The estimate here is for a map of positive associative dg algebras,
perfect outer factors and the ranks of three specified homology maps;
it identifies an explicit product that survives. These distinctions do
not settle priority. The companion script tests the chain map with
signed cell differentials, truncation and a nontrivial quotient.
\end{document}
